\documentclass{article}

\usepackage{../../Style}

\title{Elliptic PDEs}
\author{Shing Tak Lam}
\date{Lent 2023\footnote{Based on lectures by Greg Taujunskas (lectures 1 - 12) and Neshan Wickramasekera (lectures 13 - 24). Last updated \today.}}

\makeindex

\begin{document}

\maketitle

\tableofcontents

\lec

\setcounter{section}{-1}

\section{Introduction}

We will study second order elliptic PDEs on (a domain in) \(\R^n\). For example, arising from variational problems. Ultimately, we are interested in non-linear PDEs. To do this, we will first understand the linear theory.

\textbf{Setup:} Consider a domain \(\Omega \sse \R^n\), i.e. open, bounded and connected, and a function
\begin{align*}
    F : \Omega \times \R \times \R^n &\to \R \\
    (x, z, p) &\mapsto F(x, z, p)
\end{align*}
and consider the functional
\[\mcF[u] = \int_\Omega F(x, u(x), \partial u(x))\dd x\]
Note we will use \(\partial, D, \grad\) essentially interchangably. Assume \(F\) is sufficiently regular. Let \(u \in \mcS\), a suitable vector space of functions \(u : \Omega \to \R\). Frequently, \(\mcS = \HH^1(\Omega)\), or \(\mcS = C^{1, \alpha}(\Omega)\).

Suppose \(u\) minimises \(\mcF\), subject to \(u \vert_{\partial\Omega} = g\), for some given \(g : \partial\Omega \to \R\)\footnote{Boundary conditions are needed for well-posedness}. So for all \(\varphi \in \mcS\),
\[\mcF[u + t\varphi] \ge \mcF[u]\]
In particular, this means that
\[\dv{t}\bigg\vert_{t=0}\mcF[u + t\varphi] = 0\]
Or another words,
\[\dv{t}\bigg\vert_{t=0}\int_\Omega F(x, u + t \varphi, \partial u + t\partial\varphi) \dd x = 0\]
Assume enough regularity so that we can exchange the derivative and integral, we get that
\begin{equation}
    \label{eq:sec-0-calc-of-var}
    \int_\Omega(\partial_zF)(x, u, \partial u) \varphi + \partial_i\varphi(\partial p_iF)(x, u, \partial u)\dd x = 0
\end{equation}
As usual, we will use the summation convention. To ensure that \(u + t\varphi\) satisfies the correct boundary conditions, \(\varphi\vert_{\partial\Omega} = 0\). Integrate \cref{eq:sec-0-calc-of-var} by parts, we get
\[\int_{\Omega}\varphi(x)\left(\partial_z F - \partial_i\partial p_i \mcF\right)(x, u, \partial u)\dd x = 0\]
This is true for all \(\varphi \in \mcS\), and so by the fundamental lemma of calculus of variations,
\[\pdv{F}{z} - \pdv{x_i}\left(\pdv{F}{p_i}\right) = 0\]
which is the Euler-Lagrange equation for \(F\). We can rewrite this as
\begin{equation}
    \label{eq:sec-0-el-quasilinear}
    \pdv{F}{z} - \partial_i\partial_j u \pdv[2]{F}{p_i}{p_j} = 0
\end{equation}
This is now a second order quasilinear PDE in \(u\). More generally, consider
\begin{equation}
    \label{eq:sec-0-quasilinear-pde}
    a^{ij}(x, u, \partial u)\partial^2_{ij}u - b(x, u, \partial u) = 0
\end{equation}

\begin{definition}
    We say that \cref{eq:sec-0-quasilinear-pde} is \emph{elliptic} in \(\Omega\) if \(a^{ij}(x, u, \partial u)\) is a positive definite matrix in \(\Omega\).
\end{definition}

In the case of \cref{eq:sec-0-el-quasilinear}, this is equivalent to \(F\) being convex in the variable \(p\).

\begin{example}
    [Dirichlet energy] When
    \[F(x, z, p) = \abs{p}^2\]
    One gets
    \[\Delta u = 0\]
    Extremisers of this are called \emph{harmonic functions}.
\end{example}

\begin{example}
    [Minimal surfaces] When
    \[F(x, z, p) = \sqrt{1 + \abs{p}^2}\]
    We leave as an exercise to interpret \(\mcF[u]\). In this case, we get the \emph{minimal surface equation}, which is
    \begin{equation}
        \label{eq:sec-0-minimal-surface}
        \grad\vdot\left(\frac{\grad u}{\sqrt{1 + \abs{\grad u}^2}}\right) = 0
    \end{equation}
    Locally, \(\grad u\) is a constant, and so \cref{eq:sec-0-minimal-surface} looks like \(\Delta u = 0\), and so solutions have similar local properties. But the existence theory for Laplace's equation is `trivial', whereas the existence theory for \cref{eq:sec-0-minimal-surface} may fail. That is, global properties are important.

    For entire solutions (i.e. solutions defined all of \(\R^n\)), global behaviour is very different.
\end{example}

\begin{theorem}
    [Liouville] If \(u : \R^n \to \R\) is \(C^2\), \(\Delta u = 0\) and \(u\) is bounded, then \(u\) is a constant.
\end{theorem}

\begin{theorem}
    [Bernstein] (The only entire solutions to \cref{eq:sec-0-minimal-surface} in \(\R^n\) are planar (i.e. \(u\) is linear)) if and only if \(n \le 7\).
\end{theorem}

\section{Harmonic functions}

\subsection{Basic properties}

Let \(\Omega \sse \R^n\) be a domain, i.e. open and connected.

\begin{definition}
    [hamonic, subharmonic, superharmonic]
    A function \(u \in C^2(\Omega)\) is \emph{harmonic} if \(\Delta u = 0\). It is \emph{subharmonic} if \(\Delta u \ge 0\) and \emph{superharmonic} if \(\Delta u \le 0\).
\end{definition}

Let \(B_\rho(y)\) denote the open ball with centre \(y\) and radius \(\rho\). Then
\begin{theorem}
    [mean value property] If \(u \in C^2(\Omega)\) is subharmonic, and \(B_r(y) \sse \Omega\), then
    \begin{equation}
        \label{eq:mean-value-ball}
        u(y) \le \frac{1}{\omega_n r^n}\int_{B_r(y)}u(x) \dd x
    \end{equation}
    where \(\omega_n = \abs{B_1(0)}\). Moreover,
    \begin{equation}
        \label{eq:mean-value-sphere}
        u(y) \le \frac{1}{n\omega_n r^{n-1}}\int_{\partial B_r(y)}u(x)\dd x
    \end{equation}
    If \(u\) is superharmonic, then the inequalities are reversed. If \(u\) is harmonic, then equality holds.
\end{theorem}

\begin{proof}
    We have that
    \begin{align*}
        0 &\le \int_{B_\rho(y)}\Delta u \dd x \\
        \text{integrating by parts } &= \int_{\partial B_\rho(y)}\grad u \vdot w \dd x \text{ where }w\text{ is the outwards normal} \\
        &= \rho^{n-1}\int_{S^{n-1}}w\vdot \grad u(y + \rho w)\dd w \\
        &= \rho^{n-1}\int_{S^{n-1}}\pdv{\rho}u(y + \rho w)\dd w \\
    \end{align*}
    where we use the fact that \(\rho w = x - y\). Exchanging integrals and derivatives,
    \[0 \le \pdv{\rho}\int_{S^{n-1}}u(y + \rho w)\dd w\]
    Thus, the map
    \[\rho \mapsto \int_{S^{n-1}}u(y + \rho w)\dd w\]
    is increasing. Thus,
    \[\int_{S^{n-1}}u(y + \rho w)\dd w \le \int_{S^{n-1}}y(y + r w)\dd w\]
    for \(0 \le \rho \le r\). Taking the limit as \(\rho \to 0\), we get \cref{eq:mean-value-sphere}. Integrating in \(r\) to get \cref{eq:mean-value-ball}. The superharmonic case is similar. The harmonic case follows from the subharmonic and superharmonic cases.
\end{proof}

\begin{remark}
    The mean value property characterises harmonic functions. See examples sheet 1.
\end{remark}

\lec

\begin{theorem}
    [strong maximum principle] Suppose \(u \in C^2(\Omega)\) is subharmonic on \(\Omega\). Suppose there exists \(y_0 \in \Omega\) such that
    \[u(y_0) = \sup_\Omega u\]
    Then \(u\) is constant.
\end{theorem}

\begin{remark}
    If \(u\) is superharmonic, then we have a corresponding statement for when \(u(y_0) = \inf_\Omega u\). If \(u\) is harmonic, then either \(\sup\) or \(\inf\) work.
\end{remark}

\begin{proof}
    Let \(M = \sup_\Omega u < \infty\), and
    \[\Sigma = \{y \in \Omega \mid u(y) = M\}\]
    By assumption, \(\Sigma\) is non-empty, as \(y_0 \in \Sigma\). Since \(u\) is continuous, \(\Sigma\) is closed. Since \(\Omega\) is connected, it suffices to show that \(\Sigma\) is open, since this implies \(\Sigma = \Omega\).

    Pick \(y \in \Sigma\). By the mean value property, for \(\rho > 0\) such that \(\overline{B_\rho(y)} \sse \Omega\), we have that
    \[M = u(y) \le \frac{1}{\omega_n\rho^n}\int_{B_\rho(y)}u(x)\dd x\]
    So
    \[\int_{B_\rho(y)}(M - u(x))\dd x \le 0\]
    But \(M - u(x) \ge 0\), and so it must be identically zero, i.e. \(u \equiv M\) on \(B_\rho(y)\). Hence \(B_\rho(y) \sse \Sigma\), and \(\Sigma\) is open.
\end{proof}

Here, the strong maximum principle is easy as we have the mean value property. For more general elliptic PDEs, this is not the case. We will prove a weaker statement which does generalise.

\begin{theorem}
    [weak maximum principle] Suppose \(\Omega \sse \R^n\) is a bounded domain, and \(u \in C^2(\Omega) \cap C^0(\overline{\Omega})\). If \(u\) is subharmonic on \(\Omega\), then
    \[\sup_\Omega u = \sup_{\partial\Omega} u\]
\end{theorem}

\begin{proof}
    This follows from the strong maximum principle. Since \(\Omega\) is bounded, \(\sup_\Omega u\) is attained in \(\overline\Omega\). So by the maximum principle, these cannot be attained in \(\Omega\) unless \(u\) is constant.
\end{proof}

\begin{remark}
    If \(u\) is superharmonic, we replace \(\sup\) with \(\inf\), and if \(u\) is harmonic then both hold.
\end{remark}

The mean value property states that \(u\) always an average of itself. In particular, this suggest that \(u\) cannot vary too much. Can we use this to relate \(\sup u\) and \(\inf u\)?

\begin{theorem}
    [Harnack's inequality] Suppose \(u \in C^2(\Omega)\), \(u \ge 0\) and \(\Delta u = 0\) in \(\Omega\). Then if \(\Omega' \Subset \Omega\) is any bounded subdomain, we have
    \[\sup_{\Omega'}u \le C \inf_{\Omega'}u\]
    where \(C = C(n, \Omega, \Omega')\) does not depend on \(u\).
\end{theorem}

\begin{proof}
    First, choose \(y \in \Omega\) and \(\rho > 0\), such that \(\overline{B_{4\rho}(y)} \sse \Omega\). Choose \(x_1, x_2 \in B_\rho(y)\). By the mean value property,
    \[u(x_1) = \frac{1}{\omega_n\rho^n}\int_{B_\rho(x_1)}u \le \frac{1}{\omega_n\rho^n}\int_{B_{2\rho}(y)}u\]
    On the other hand,
    \[u(x_2) = \frac{1}{\omega_n(3\rho)^n}\int_{B_{3\rho}(y)}u \ge \frac{1}{\omega_n(3\rho)^n}\int_{B_{2\rho(y)}}u\]
    Combining these, we see that
    \[u(x_1) \le 3^n u(x_2)\]
    for all \(x_1, x_2 \in B_\rho(y)\). So Harnack holds locally in balls, with constant independent of \(u\). It is also independent of \(\rho, y\) as long as \(\rho\) is sufficiently small. Now choose \(x_1, x_2\) in \(\overline{\Omega'} \sse \Omega\), such that
    \[\sup_{\Omega'}u = u(x_1) \qqtext{and} \inf_{\Omega'}u = u(x_2)\]
    By path connectedness of \(\Omega\), there exists a continuous map \(\gamma : [0, 1] \sse \overline{\Omega'}\), with \(\gamma(0) = x_1\) and \(\gamma(1) = x_2\). Choose \(\rho\) such that \(4\rho < d(\gamma, \partial\Omega)\), and \(N = N(\Omega', \Omega)\) such that we can cover \(\gamma\) by \(N\)-balls of radius \(\rho\)\footnote{This follows from the fact that \(\Omega'\) is relatively compact}.

    Apply the local result to each ball, and we get
    \[u(x_1) \le (3^n)^N u(x_2) = 3^{nN}u(x_2)\]
\end{proof}

\begin{theorem}
    [derivative estimates] Suppose \(u \in C^3(\Omega)\) is harmonic on \(\Omega\). Then if \(\overline{B_\rho(y)} \sse \Omega\), we have that
    \[\abs{\DD u(y)} \le \frac{C}{\rho}\sup_{\partial B_\rho(y)}\abs{u}\]
    where \(C = C(n)\).
\end{theorem}

\begin{proof}
    Since \(\Delta u = 0\),
    \[0 = \DD_i(\Delta u) = \Delta(\DD_i u)\]
    So \(\DD_i u\) is harmonic. By the mean value property,
    \begin{align*}
        \DD_i u(y) &= \frac{1}{\omega_n\rho^n}\int_{B_\rho(y)}u \\
        &= \frac{1}{\omega_n\rho^n}\int_{B_\rho(y)}\grad\vdot(0, \dots, u, \dots, 0) \dd x\\
        &= \frac{1}{\omega_n\rho^n}\int_{\partial B_\rho(y)}u(x)\vdot \nu_i(x)\dd x
    \end{align*}
    where \(\nu(x)\) is the unit normal at \(x\). But \(\abs{\nu_i(x)} \le 1\), and so
    \begin{align*}
        \abs{\DD_i u(y)} \le \frac{1}{\omega_n \rho^n}\sup_{\partial B_\rho(y)}\abs{u}\int_{\partial B_\rho(y)}\dd x = \frac{n}{\rho}\sup_{\partial B_\rho(y)}\abs{u}
    \end{align*}
\end{proof}

\begin{remark}
    We can apply this result repeatedly, to get that for \(\Omega' \sse \Omega'' \sse \Omega\), and any multi-index \(\alpha\), if \(u \in C^{\abs{\alpha} + 2}(\Omega)\), with \(\Delta u = 0\) in \(\Omega\), then
    \[\sup_{\Omega''}\abs{\DD^\alpha u} \le C \sup_{\Omega''}\abs{u}\]
    for some \(C = C(n, \alpha, \Omega, \Omega')\).

    That is,
    \[\norm{\DD^\alpha u}_{L^\infty(\Omega')} \le C\norm{u}_{L^\infty(\Omega')}\]
    By the mean value property, for some \(y \in \overline{\Omega''} \sse \Omega\),
    \[\sup_{\Omega''}u = \abs{u(y)} = \abs{c\int_{B_\rho(y)}u(x)\dd x} \le C\int_\Omega \abs{u}\]
    So
    \[\norm{\DD^\alpha(u)}_{L^\infty(\Omega')} \le C\norm{u}_{L^1(\Omega)}\]
\end{remark}

\begin{theorem}
    [uniqueness of solutions for the Dirichlet problem] Suppose \(\Omega\) is bounded, and \(u_1, u_2 \in C^2(\Omega) \cap C^0(\overline{\Omega})\) with
    \[\begin{cases}
        \Delta u_1 = \Delta u_2 & \text{in }\Omega \\
        u_1 = u_2 &\text{on }\partial\Omega
    \end{cases}\] 
    Then \(u_1 = u_2\) in \(\overline\Omega\).
\end{theorem}

\begin{proof}
    Set \(w = u_2 - u_2\). Then \(w\) is harmonic function in \(\Omega\), and \(w = 0\) on \(\partial\Omega\). By the weak maximum principle, we get that \(w = 0\) in \(\Omega\).
\end{proof}

\begin{remark}
    Of course, we can integrate by parts to get the same result, but the weak maximum principle will apply for non divergence form equations.
\end{remark}

\lec

\begin{theorem}
    [Liouville] If \(u \in C^\infty(\R^n)\) is harmonic, and grows sublinearly at infinity. Then \(u\) is constant.
\end{theorem}

\begin{remark}
    ``growing sublinearly'' means that
    \[\abs{u(x)} \le C\left(1 + \abs{x}^\alpha\right)\]
    where \(0 < \alpha < 1\).
\end{remark}

\begin{proof}
    From derivative estimates, we know that for all \(x \in \R^n\), we have that
    \[\abs{\DD u(y)} \le \frac{C}{\rho}\sup_{B_\rho(y)}\abs{u}\]
    Plugging in the growth assumption, we get that
    \[\abs{\DD u(y)} \le \frac{C}{\rho}(1 + (\rho + \abs{y})^\alpha)\]
    Taking \(\rho \to \infty\), we get that \(\DD u(y) = 0\). But \(y\) was arbitrary, and so we are done.
\end{proof}

\subsection{Existence theory for harmonic functions}

We will consider the classical problem of solving the Dirichlet problem for the Laplacian on bounded domains \(\Omega \sse \R^n\), and \(\varphi : \overline\Omega \to \R\) continuous. We wish to find \(u \in C^\infty(\Omega) \cap C^0(\overline\Omega)\) such that
\[\begin{cases}
    \Delta u = 0 & \text{in }\Omega \\
    u = \varphi & \text{on }\partial\Omega
\end{cases}\]
We will assume for simplicity that \(\partial\Omega\) is smooth, and \(\varphi \in C^\infty(\Omega)\).

We have (at least) three methods to solve this problem.
\begin{enumerate}
    \item Hilbert space methods (c.f. Analysis of PDEs). Use the Riesz representation theorem to obtain a solution \(u \in \HH^1(\Omega)\), and deal with regularity afterwards. This relies on the equation being linear.
    \item Direct method of calculus of variations. We can rephrase \(\Delta u = 0\) as a variational problem. That is, the Euler-Lagrange equation of
    \[\int\abs{\DD u}^2\]
    and prove existence using the functional.
    \item Perron's method. Use the fact that solvability in balls impies solvability in more general domains. This method is based on maximum principles.
\end{enumerate}

\begin{remark}
    In all cases, we obtain a rougher solution first, and improve regulaity later.
\end{remark}

We will focus on the second method. Define
\[\msS = \left\{w \in \HH^1(\Omega) \mid w - \varphi \in \HH_0^1(\Omega)\right\}\]
That is, \(\HH^1\) functions which agree with \(\varphi\) on the boundary. Clearly \(\varphi \in \msS\), and so \(\msS\) is non-empty. Set
\[\mcE[w] = \int_\Omega \abs{\DD w}^2\]
and define
\[\beta = \inf_{w \in \msS} \mcE[w]\]
By the definition of the infimum, there exists a seqeunce \((w_j) \sse \msS\), such that
\[\mcE[w_j] \to \beta\]
We want to extract a convergent subsequence and show that its limit is a solution. Clearly for \(j\) large, we have that
\[\int_\Omega \abs{\DD w_j}^2 \le \beta + 1\]
Since \(w_j - \varphi \in \HH^1_0(\Omega)\), by the Poincar\'e inequality,
\[\int_\Omega\abs{w_j - \varphi}^2 \le C\int_\Omega \abs{\DD(w_j - \varphi)}^2\]
In particular, this implies that
\[\norm{w_j}_{L^2(\Omega)}^2 \le C(\Omega, \varphi, \beta) < \infty\]
Indeed,
\[\norm{w_j}_{L^2(\Omega)}^2 \le C\norm{\DD w_j - \DD \varphi}_{L^2(\Omega)}^2 \le C(\Omega, \varphi, \beta)\]
Expanding the left hand side, we get
\[\norm{w_j}_{L^2(\Omega)}^2 - 2 \inner{w_j, \varphi}_{L^2(\Omega)} \le C(\Omega, \varphi, \beta)\]
In particular, by Young's inequality,
\[\norm{w_j}_{L^2(\Omega)}^2 \le C(\varphi, \Omega, \beta) + \epsilon\norm{w_j}_{L^2(\Omega)}^2 + \frac{1}{\epsilon}\norm{\varphi}^2_{L^2(\Omega)}\]
Take \(0 < \epsilon < 1\), and rearrange.

So we have that \(\norm{w_j}_{\HH^1(\Omega)}^2 \le C\), so by Banach-Alaoglu,
\[w_{j_k} \weak w \text{ in }\HH^1(\Omega)\]
By Rellich-Kondrachov\footnote{Recall \begin{theorem}
    [Rellich-Kondrachov] Let \(\Omega\) be bounded, \(1 \le p < n\), then
    \[W^{1, p}(\Omega) \inj L^{p^*}(\Omega)\]
    and
    \[W^{1, p}(\Omega) \Subset L^q(\Omega)\]
    for \(1 \le q < p^*\). Here, \(p^*\) is the conjugate
    \[p^* = \frac{np}{n - p}\]
\end{theorem}
In particular, when \(p = 2\), \(p^* > 2\) if \(n > 2\).},
\[w_{j_k} \to w \text{ in }L^2(\Omega)\]
for some \(w \in \HH^1(\Omega)\). Hence for all \(v \in \HH^1(\Omega)\), we have that
\[\int_{\Omega}\DD w_j \vdot \DD v \to \int_\Omega \DD w \vdot \DD v\]
Also, clearly we have that \(w_{j_k} - \varphi \weak w - \varphi\) in \(\HH^1(\Omega)\), as \(\varphi\) is smooth. But \(w_{j_k} - \varphi \in \HH^1_0(\Omega)\), but \(\HH^1_0(\Omega)\sse \HH^1(\Omega)\) is norm closed, so it is weakly closed\footnote{This follows from Hahn-Banach. In fact, it follows for any convex subset of a Banach space.

\begin{lemma}
    Let \(X\) be a Banach space, for \(C \sse X\) convex, \(C\) is norm-closed if and only if \(C\) is weakly closed.
\end{lemma}
\begin{proof}
    Weak closed implies norm closed is clear. For the converse, we can show that \(X \setminus C\) is weakly open. Let \(x_0 \in X\setminus C\), by the Hahn-Banach separation, there exists \(\phi \in X'\) such that
    \[\varphi\vert_{C} = 0 \qqtext{and}\phi(x_0) \ne 0\]
    Then
    \[\left\{x \in X \ \bigg\vert\ \abs{\phi(x)} > \frac{1}{2}\abs{\phi(x_0)}\right\} \sse X \setminus C\]
    is a weakly open neighbourhood of \(x_0\).
\end{proof}}.
Hence \(w - \varphi \in \HH^1_0(\Omega)\), so \(w \in \msS\). Finally, since \(\mcE[\cdot]\) is \emph{sequentially weakly lower semicontinuous}\footnote{
    That is, if \(u_j \weak u\) in \(\HH^1(\Omega)\), then
    \[\mcE[u] \le \liminf \mcE[u_j]\]
    To see this, note that
    \[\int_\Omega\DD u_j \vdot \DD v \to \int_\Omega \DD u \vdot \DD v\]
    Setting \(u = v\), we see that
    \[\int_\Omega\DD u_j \vdot \DD u \to \int_\Omega \abs{\DD u}^2\]
    Thus,
    \begin{align*}
        \mcE[u] &= \lim\int_\Omega \DD u_j \vdot \DD u \\
        &= \liminf \int_\Omega \DD u_j \vdot \DD u \\
        &\le \liminf \mcE[u_j]^{1/2}\mcE[u]^{1/2}
    \end{align*}
} in \(\HH^1(\Omega)\), we have that
\[\mcE[w] \le \liminf \mcE[w_{j_k}] = \beta\]
Hence \(\mcE[w] = \beta\). We have found a global minimum \(w\). That is, for all \(v \in \HH^1_0(\Omega)\), \(w + tv \in \msS\),
\[f(t) = \mcE[w + tv] \ge \mcE[w]\]
Thus, \(f'(0) = 0\). But
\begin{align*}
    f'(0) &= \DD \mcE[w](v) \\
    &=\lim_{t \to 0}\frac{\mcE[w + t v] - \mcE[w]}{t} \\
    &= 2\int_\Omega \DD w \vdot \DD v
\end{align*}
In particular,
\[\int_\Omega \DD w \vdot \DD v = 0\]
is the weak formulation of \(\Delta w = 0\). We will next upgrade the regularity, and show that weak solutions are in fact smooth.

\subsection{Interior regularity}

We wish to improve the regularity of the weak solution. What we have shown is that there exists a weak solution \(u \in L^1(\Omega)\), such that
\[\int_\Omega u \Delta v = 0\]
for all \(v \in C_c^\infty(\Omega)\).

\begin{theorem}
    [Weyl's lemma] Weakly harmonic functions are smooth. That is, for \(\Omega \sse \R^n\) a domain, \(u \in L^1_\loc(\Omega)\), if we have \(u\) is a weak solution to Laplace's equation, then \(u\) is \(C^\infty\) and \(\Delta u = 0\) in \(\Omega\).
\end{theorem}

\begin{proof}
    Mollify \(u\). Take \(\varphi \in C^\infty(\R^n)\), such that
    \begin{itemize}
        \item \(0 \le \varphi\),
        \item \(\varphi(x) = 0\) for \(\abs{x} \ge 1\),
        \item \(\int_{\R^n}\varphi = 1\),
        \item \(\varphi\) is radially symmetric\footnote{This is not a standard assumption for mollifiers, but we can assume this.}.
    \end{itemize}
    For \(\sigma > 0\), set
    \[\varphi_\sigma(x) = \frac{1}{\sigma^n}\varphi\left(\frac{x}{\sigma}\right)\]
    Then \(\varphi_\sigma \in C_c^\infty(B_\sigma(0))\) is nonnegative and has integral \(1\).

    Define
    \[u_\sigma(x) = (\varphi_\sigma * u)(x)\]
    This is well defined for
    \[x \in \Omega_\sigma = \{x \in \Omega \mid d(x, \partial\Omega) \ge \sigma\}\]
    Then \(u_\sigma\) is smooth, \(u_\sigma \to u\) in \(L^1_\loc(\Omega)\). Moreover, \(\Delta u_\sigma = 0\). To see this,
    \[\pdv{x^i}u_\sigma = \int_\Omega u(y)\pdv{x^i}\varphi_\sigma(x - y)\dd y = -\int_\Omega u(y)\pdv{y^i}\varphi_\sigma(x - y)\dd y\]
    and so
    \[\Delta_x u_\sigma(x) = \int_\Omega u(y)\Delta_y \varphi_\sigma(x - y)\dd y = 0\]
    as \(u\) is weakly harmonic.

    By the a priori derivative estimates for harmonic functions, for \(\Omega' \Subset \Omega\),
    \[\sup_{\Omega'}\abs{\DD^\alpha u_\sigma} \le C \int_{\Omega'_{\sigma_1}}\abs{u_\sigma}\]
    for some \(\sigma_1(\Omega')\) small, where
    \[\Omega'_\sigma = \Omega' \cup \{x \in \Omega \mid d(x, \partial\Omega') < \sigma\}\]
    Since \(u_\sigma \to u\) in \(L^1_\loc(\Omega)\), for \(\sigma\) small enough,
    \[C \int_{\Omega'_{\sigma_1}}\abs{u_\sigma} \le C\int_{\Omega'_{\sigma'}}(\abs{u} + 1)\]
    Hence
    \[\sup_{\Omega'}\abs{\DD^\alpha u_\sigma} \le C\int_{\Omega'_{\sigma_1}}(\abs{u} + 1)\]
    So \(\DD^\alpha u_\sigma\) is uniformly bounded in \(L^\infty(\Omega')\). Hence (as bounded derivatives imply equicontinuity) by Arzela-Ascoli, there exists a subsequence \((\sigma_j)\) such that \(\sigma_j \to 0\), and there exists \(\widetilde u \in C^\infty(\Omega)\) such that \(u_{\sigma_j} \to \widetilde u\) in \(C^k(\Omega')\) for all \(k\). Hence
    \[\Delta \widetilde u = \lim_{j \to \infty} \Delta u_{\sigma_j} = 0\]
    in \(\Omega\), as \(\Omega'\) was arbitrary. By properties of mollifiers, \(u_\sigma \to u\) a.e. in \(\Omega\), and so \(\widetilde u = u\) a.e.
\end{proof}

\begin{remark}
    We do not say anything about boundary regularity. It is possible to get (at least) \(u \in C^0(\overline\Omega)\).
\end{remark}

Let's now improve our \(C^\infty\) existence result to \(C^0\).

\begin{theorem}
    [existence and uniqueness for the Dirichlet problem with \(C^0\) data] Suppose \(\Omega\) is bounded with \(\partial\Omega\) sufficiently regular, then for any \(\varphi \in C^0(\partial\Omega)\), there exists a unique \(u \in C^\infty(\Omega) \cap C^0(\overline\Omega)\) solving
    \[\begin{cases}
        \Delta u = 0 &\text{in }\Omega \\
        u = \varphi&\text{on }\partial\Omega
    \end{cases}\]
\end{theorem}

\begin{remark}
    We might have 
    \[\int_\Omega\abs{\DD u}^2 = \infty\]
    for this solution.
\end{remark}

\begin{proof}
    Choose a sequence \((\varphi_n) \sse C^\infty(\R^n)\), such that \(\varphi_n \to \varphi\) on \(\partial \Omega\), in \(L^\infty\). Then we know that there exists \(u_n \in C^\infty(\Omega) \cap C^0(\overline\Omega)\), such that
    \[\begin{cases}
        \Delta u_n = 0 &\text{in }\Omega \\
        u_n = \varphi_n &\text{on }\partial \Omega
    \end{cases}\]
    Then for all \(n, m \in \N\), \(\Delta(u_n - u_m) = 0\) in \(\Omega\), and \(u_n - u_m = \varphi_n - \varphi_m\) on \(\partial\Omega\). By the weak maximum principle,
    \[\sup_{\overline\Omega}\abs{u_n - u_m} \le \sup_{\partial\Omega}\abs{u_n - u_m} = \norm{\varphi_n - \varphi_m}_{L^\infty} \to 0\]
    as \(n, m \to \infty\). So \((u_n)\) is Cauchy in \(C^0(\overline\Omega)\), which is a Banach space, so by completeness, there exists \(u \in C^0(\overline\Omega)\) such that \(u_n \to u\) uniformly on \(\overline\Omega\). In particular, \(u = \varphi\) on \(\partial\Omega\).

    By the derivative estimates, \((u_n)\) converges in \(C^k(\Omega')\), for any \(\Omega' \Subset \Omega\), and so \(u\) is smooth in the interior. Thus, \(\Delta u = 0\) in the interior.
\end{proof}

\begin{remark}
    A sufficient condition for regularity is that \(\partial\Omega\) is \(C^2\). More generally, it is enough to have the \emph{exterior sphere condition}, which says that for all \(z \in \partial\Omega\), there exists \(B_\rho(y)\), such that \(\overline{B_\rho(y)} \cap \partial\Omega = \{z\}\). There exists bounded domains in which this fails, and the conclusion of the theorem fails in that case as well. For example, when the boundary has a cusp.
\end{remark}

\section{General second order elliptic operators}

From now on, write
\[Lu = a^{ij}\partial_i\partial_j u + b^i\partial_i u + cu\]
We will work on \(\Omega \sse \R^n\) open, \(u \in C^2(\Omega)\), \(a^{ij}, b^i, c : \Omega \to \R\) and consider the Dirichlet problem
\[\begin{cases}
    Lu = f &\text{in }\Omega \\
    u = \varphi&\text{on }\partial\Omega
\end{cases}\]
for given \(f\) and \(\varphi\). If we can write \(L\) in divergence form
\[Lu = \partial_i(a^{ij}\partial_ju) + \widehat{b}^i \partial_i u + cu\]
then we can use Hilbert space methods, as in Analysis of PDEs. But if \(a^{ij}\) is only \(C^0\) say, we will need Schauder theory.

The idea is to deform \(L\) into \(\Delta\) using a series of rescalings. In particular, this does not involve Sobolev spaces. Since \(u \in C^2(\Omega)\), we can assume \(a^{ij}\) is symmetric.

\begin{definition}
    We say that \(L\) is \emph{elliptic} in \(\Omega\) if the matrix \((a^{ij})\) is positive definite in \(\Omega\). So that
    \[0 \le \lambda(x)\abs{\xi}^2 \le a^{ij}(x)\xi^i\xi^j \le \Lambda(x)\abs{\xi^2}\]
    for all \(\xi \in \R^n\) non-zero. In particular, \(\lambda\) is the minimum eigenvalue, and \(\Lambda\) is the maximum.

    \(L\) is \emph{strictly elliptic} if there exists \(\lambda_0\) such that \(0 < \lambda_0 \le \lambda(x)\) for all \(x\). \(L\) is \emph{uniformly elliptic} if it is elliptic, and \(\Lambda(x)/\lambda(x)\) is uniformly bounded.
\end{definition}

\begin{remark}
    In general, uniformly elliptic does not imply strictly elliptic.
\end{remark}

\begin{example}
    The minimal surface equation
    \[\grad\vdot\left(\frac{\DD u}{\sqrt{1 + \abs{\DD u}^2}}\right) = 0\]
    has
    \[a^{ij} = \left(\delta_{ij} - \frac{\DD_i u \DD_j u}{1+ \abs{\DD u}^2}\right)\frac{1}{\sqrt{1 + \abs{\DD u}^2}}\]
    This is elliptic but not uniformly elliptic.
\end{example}

\lec

We are interested in general second order elliptic operators, with \(a^{ij}, b^i, c \in C^{0, \alpha}(\Omega)\). In particular, we cannot write them in divergence form, as \(a^{ij}\) are not \(C^1\).

\subsection{Basic properties}

\begin{theorem}
    [weak maximum principle] Suppose that \(L\) is elliptic and that
    \[\sup_\Omega \abs{\frac{b^i}{\lambda}} < \infty \tag{\((*)\)}\]
    for some \(i\). Moreover, suppose \(\Omega\) is bounded, open, \(u \in C^2(\Omega) \cap C^0(\overline\Omega)\), such that \(Lu \ge 0\)\footnote{We say that \(u\) is a \emph{subsolution}.}.

    Then
    \begin{itemize}
        \item if \(c = 0\), then
        \[\sup_\Omega u = \sup_{\partial\Omega} u\]
        \item if \(c \le 0\), then
        \[\sup_\Omega u \le \sup_{\partial\Omega}u^+\]
        where \(u^+ = \max(u, 0)\).
    \end{itemize}
\end{theorem}

\begin{remark}
    The assumption that \(c \le 0\) is crucial. For example, when \(n = 1\), let \(\Omega = (0, \pi)\), and \(u'' + u = 0\). Then \(u(x) = \sin(x)\) is a solution. But then \(\sup_\Omega u = 1\), and \(\sup_{\partial\Omega}u^+ = 0\).

    For \(n = 2\), \(\Omega = (0, \pi)^2\), and \(\Delta u + 2u = 0\). Then
    \[u(x, y) = \sin(x)\sin(y)\]
    has the same properties as the above.
\end{remark}

\begin{proof}
    First suppose \(c = 0\). If \(Lu > 0\) in \(\Omega\), then in fact the strong maximum principle holds. Indeed, if \(x_0 \in \Omega\) is a local maximum, then
    \[\partial_i u(x_0) = 0 \qqtext{and} \partial_i\partial_j u(x_0) \preccurlyeq 0\]
    Since \(a^{ij}(x_0) \succcurlyeq 0\), we have that 
    \[a^{ij}\partial_i\partial_ju(x_0) = \tr(A \grad^2u(x_0)) \le 0\]
    Hence \(0 < Lu(x_0) = a^{ij}\partial_i\partial_j u(x_0) + b^i\partial_iu(x_0) \le 0\). Contradiction. More generally, if \(Lu \ge 0\) in \(\Omega\), consider
    \[v(x) = e^{\gamma x_1}\]
    for some \(\gamma > 0\) to be chosen. Here, we assume without loss of generality that \((*)\) holds for \(i = 1\). Then
    \[\partial_1 v = \gamma e^{\gamma x_1} \qqtext{and}\partial_i v = 0 \text{ for }i \ne 1\]
    and
    \[\partial_1\partial_1 v \gamma^2 e^{\gamma x_1}\]
    and all other second derivatives are zero. So
    \begin{align*}
        Lv &= e^{\gamma x_1}(a^{11}\gamma^2 + b^1\gamma) \\
        &\ge e^{\gamma x_1}(\lambda \gamma^2 + b^1\gamma) \\
        &= \gamma e^{\gamma x_1}\left(\gamma^2 + \frac{b^1}{\lambda}\gamma\right)
    \end{align*}
    This is positive for \(\gamma\) large enough. Hence
    \[L(u + \epsilon v) > 0\]
    for all \(\epsilon > 0\). By the above,
    \begin{align*}
        u(x) &\le \sup_{\Omega}(u + \epsilon v) \quad\text{ as }v \ge 0\\
        &\le \sup_{\partial\Omega}(u + \epsilon v) \\
        &\le \sup_{\partial\Omega} u + \epsilon\sup_{\partial\Omega}v
    \end{align*}
    Taking \(\epsilon \to 0\), we get that \(u(x) \le \sup_{\partial\Omega} u\). Since this is true for all \(x \in \Omega\), then
    \[\sup_\Omega u \le \sup_{\partial \Omega} u\]
    The reverse inequality is trivial.

    Now suppose \(c \le 0\). Define
    \[L_0u = a^{ij}\partial_i\partial_j u + b^i\partial_i u\]
    and consider
    \[\Omega^+ = \{x \in \Omega \mid u(x) > 0\}\]
    Since \(cu \le 0\) on \(\Omega^+\),
    \[L_0u = Lu - cu \ge 0\]
    on \(\Omega^+\). Note if \(\Omega^+ = \emptyset\), then \(\sup_\Omega u \le 0\), and \(u^+ = 0\), and the conclusion is trivial. Thus, without loss of generality assume \(\Omega^+ \ne \emptyset\). Then there exists \(x_0\in \partial\Omega \cap \partial\Omega^+\), with \(u(x_0) \ge 0\). If not, then
    \[\partial\Omega^+  \cap \partial\Omega = \emptyset\]
    and so \(\partial\Omega^+ \sse \Omega\), and so it \(\partial\Omega^+ \sse \Omega \setminus \Omega^+\).  So \(u\vert_{\partial\Omega^+} \le 0\). But this contradicts the first part for \(L_0\) on \(\Omega^+\).

    Hence
    \begin{align*}
        \sup_\Omega u &= \sup_{\Omega^+} u \\
        &= \sup_{\partial\Omega^+} u \\
        &\le \sup_{\partial\Omega} u \\
        &\le \sup_{\partial\Omega} u^+
    \end{align*}
\end{proof}

\begin{corollary}
    Let \(\Omega\) be bounded open, \(u \in C^2(\Omega) \cap C^0(\overline\Omega)\), \(L\) is elliptic, with the same bound on \(\Omega\), and \(c \le 0\) in \(\Omega\). Then
    \begin{enumerate}
        \item if \(Lu \le 0\) in \(\Omega\), then
        \[\inf_{\Omega} u \ge \inf_{\partial\Omega} u^-\]
        where \(u^- = \min(u, 0)\).
        \item if \(Lu = 0\), then
        \[\sup_{\partial\Omega}\abs{u} = \sup_\Omega \abs{u}\]
    \end{enumerate}
\end{corollary}

\begin{proof}
    Exercise.
\end{proof}

\begin{corollary}
    Let \(L\) as above, and suppose we have \(u, v, w \in C^2(\Omega) \cap C^0(\overline\Omega)\), such that
    \[Lu\ge 0 \quad Lv = 0 \qqtext{and} Lw \le 0\]
    Then
    \begin{enumerate}
        \item if \(u \le v\) on \(\partial\Omega\), then \(u \le v\) in \(\overline{\Omega}\).
        \item if \(v \le w\) on \(\partial\Omega\), then \(v \le w\) in \(\overline{\Omega}\).
    \end{enumerate}
\end{corollary}

\begin{proof}
    Exercise.
\end{proof}

We want to build towards a strong maximum principle. For this, we need

\begin{theorem}
    [Hopf boundary point lemma] Let \(\Omega \sse \R^n\) be open, and take \(y \in \partial\Omega\). Suppose \(\partial\Omega\) satisfies the \emph{interior sphere condition}\footnote{That is, there exists \(R > 0\), \(z \in \Omega\), such that \(B_R(z) \sse \Omega\), \(y \in \partial B_R(z)\).} at \(y\). Let \(L\) be uniformly elliptic in \(\Omega\), with
    \[\sup_{\Omega}\frac{\abs{b}}{\lambda} + \sup_\Omega \abs{c}{\lambda} < \infty\]
    Suppose \(u \in C^2(\Omega) \cap C^0(\{y\} \cup \Omega)\), such that \(u(y) \ge u(x)\) for all \(x \in \Omega\), and \(Lu \ge 0\) in \(\Omega\). Finally, assume one of the following holds:
    \begin{enumerate}[(i)]
        \item \(c = 0\) in \(\Omega\),
        \item \(c \le 0\) in \(\Omega\) and \(u(y) \ge 0\),
        \item \(u(y) = 0\)
    \end{enumerate}
    Then
    \[\pdv{u}{\nu}(y) > 0\]
    if it exists, where \(\nu\) is the outwards pointing normal at \(y\) to \(\partial B_R(z)\), coming from the interior sphere condition.
\end{theorem}

\begin{remark}
    The weak maximum principle implies that
    \[\pdv{u}{\nu} \ge 0\]
    and so the content of the theorem is the strict inequality.
\end{remark}

\begin{proof}
    Let
    \[A = B_R(z) \setminus B_r(z)\]
    for some \(0 < r < R\). We will first solve cases (i) and (ii). On \(A\), consider
    \[v(x) = e^{-\alpha\abs{x - z}^2} - e^{-\alpha R^2}\]
    First note that on \(A\), \(v > 0\).
    \[\partial_i v(x) = -2\alpha(x_i - z_i)e^{-\alpha\abs{x - z}^2}\]
    and
    \[\partial_i\partial_jv(x) = -2\alpha\delta_{ij}e^{-\alpha\abs{x - z}^2} + 4\alpha^2(x_i - z_i)(x_j - z_j)e^{-\alpha\abs{x - z}^2s}\]
    So on \(A\),
    \[Lv = e^{-\alpha \abs{x - z}^2}(a^{ij}4\alpha^2(x_i - z_i)(x_j - z_j) - 2\alpha a^{ii} - 2\alpha b^i(x_i - z_i) + c) - ce^{-\alpha R^2}\]
    By ellipticity and the sign of \(c\),
    \begin{align*}
        Lv &\ge 2^{-\alpha\abs{x - z}^2}(4\alpha^2 \lambda(x)\abs{x - z}^2 - 2\alpha n \Lambda(x) - 2\alpha \abs{b}\abs{x - z} - \abs{c}) \\
        &\ge e^{-\alpha\abs{x - z}^2}\lambda(x)\left(\alpha^2 R^2 - 2\alpha n \sup_\Omega\frac{\Lambda}{\lambda} - \alpha R \sup_\Omega \frac{\abs{b}}{\lambda} - \sup_\Omega \frac{\abs{c}}{\lambda}\right)
    \end{align*}
    where \(r\) is chosen such that \(\abs{x - z}^2 \ge (R/2)^2\). In particular, this is positive for \(\alpha\) large enough. Fix such an \(\alpha\).
\lec

    Set \(w(x) = u(x) - u(y) + \epsilon v(x)\), for some small \(\epsilon\) to be determined. Now
    \[Lw = Lu + \epsilon Lv - cu(y) \ge 0\]
    in \(A\), by above (and the assumptions). Also, 
    \[v\vert_{\partial B_R(x)} = 0\]
    and because \(u(x) \le u(y)\) on \(\overline\Omega\), so
    \[w\vert_{\partial B_R(z)} \le 0\]
    Also, \(u(x) < u(y)\) on \(\partial_{B_r(z)}\), so we can choose \(\epsilon\) small enough such that
    \[w\vert_{\partial B_r(z)} < 0\]
    So \(w \vert_{\partial A} \le 0\). Apply the weak maximum principle to \(w\) in \(A\), we get \(w \le 0\), and so
    \[u(x) - u(y) + \epsilon v(x) \le 0\]
    in \(A\). Choose \(t < 0\), and we have that
    \[\frac{u(y + t\nu) - u(y)}{t} \ge -\epsilon \frac{v(y + t\nu) - v(y)}{t}\]
    Note \(v(y) = 0\). Sending \(t \to 0\),
    \begin{align*}
        {\pdv{u}{\nu}}(y) &\ge -\epsilon{\pdv{v}{\nu}}(y) \\
        &= -\epsilon\partial_i v(y)\left(\frac{y_i - z_i}{R}\right) \\
        &= 2\alpha\epsilon R e^{-\alpha R^2} \\
        & > 0
    \end{align*}

    For case (iii), consider
    \[\widetilde L = L - c^+\]
    So \(\widetilde L u = Lu - c^+ u \ge 0\), and we can apply the above to \(\widetilde L\).
\end{proof}

\begin{theorem}
    [strong maximum principle] Suppose \(\Omega \sse \R^n\) is a domain, with non-empty boundary \(\partial \Omega\) satisfying the interior sphere condition for all \(y \in \partial\Omega\). Let \(L\) be uniformly elliptic, with
    \[\sup_\Omega \left(\frac{\abs{b} + \abs{c}}{\lambda}\right) < \infty\]
    Suppose \(u \in C^2(\Omega)\), with \(M = \sup_\Omega u < \infty\), and \(Lu \ge 0\) on \(\Omega\). Suppose (at least) one of the following holds
    \begin{enumerate}[(i)]
        \item if \(c = 0\), and \(u(y) = M\) for some \(y \in \Omega\).
        \item if \(c \le 0, M \ge 0\) and \(u(y) = M\) for some \(y \in \Omega\),
        \item if \(M = 0\) and \(u(y) = M = 0\) for some \(y \in \Omega\).
    \end{enumerate}
    Then \(u\) is constant.
\end{theorem}

\begin{proof}
    Let
    \[\Sigma = \{x \in \Omega \mid u(x) = M\}\]
    By continuity, \(\Sigma\) is closed in \(\Omega\). Suppose \(\Omega \setminus \Sigma \ne \emptyset\). Choose \(z \in \Omega \setminus \Sigma\), such that
    \[d(z, \partial\Omega) > d(z, \partial\Sigma)\]
    To see this, first choose \(z_1 \in \partial\Sigma \cap \Omega\). Then choose \(\rho_1 > 0\) such that \(B_{\rho_1}(z_1) \sse \Omega\). Then choose any
    \[z \in B_{\rho_1/2}(z_1) \setminus \Sigma\]
    Let
    \[R = \sup\{\rho \mid B_\rho(z) \sse \Omega \setminus \Sigma\}\]
    By construction, there exists \(y \in \partial B_R(z) \cap \Sigma\). Since \(\DD u(y) = 0\), this contradicts the Hopf boundary point lemma. So \(\Omega \setminus \Sigma = \emptyset\), i.e. \(\Omega = \Sigma\), so \(u\) is constant.

    The assumption (i), (ii) and (iii) are so we can apply the Hopf boundary point lemma.
\end{proof}

Some corollaries of this:

\begin{corollary}
    [comparison principle]
    Let
    \[L = a^{ij}\partial_i\partial_j + b^i\partial_i + c\]
    be uniformly elliptic in \(\Omega \sse \R^n\), with
    \[\sup_{\Omega}\left(\frac{\abs{b} + \abs{c}}{\lambda}\right) < \infty\]
    Suppose \(u, v \in C^2(\Omega)\), such that \(Lu \ge Lv\) and \(u \le v\) in \(\Omega\). Then
    \begin{itemize}
        \item either \(u  =v\) on \(\Omega\),
        \item or \(u < v\) on \(\Omega\),
    \end{itemize}
\end{corollary}

\begin{proof}
    \(L(u - v) \ge 0\) in \(\Omega\), and \(u - v \le 0\). So if there exists \(x_0 \in \Omega\), with \(u(x_0) = v(x_0)\), then the strong maximum principle implies \(u = v\). If not, then \(u < v\) in \(\Omega\).
\end{proof}

\begin{corollary}
    [uniqueness for the Neumann problem] Suppose \(\Omega \sse \R^n\) is a bounded domain, and \(\partial\Omega\) satisfies the interior sphere condition at each point. Suppose \(L\) is uniformly elliptic, with
    \[\sup_{\Omega}\left(\frac{\abs{b} + \abs{c}}{\lambda}\right) < \infty\]
    and \(c \le 0\). Then if \(u_1, u_2 \in C^2(\Omega) \cap C^0(\overline\Omega)\) is such that
    \begin{align*}
        Lu_i = f &\text{ in }\Omega \\
        \pdv{u_i}{\nu} = g &\text{ on }\partial\Omega
    \end{align*}
    for some \(f : \Omega \to \R, g : \partial\Omega \to \R\). Then \(u_1 - u_2\) is constant.
\end{corollary}

\begin{proof}
    Let \(u = u_1 - u_2\). This satisfies the Neumann problem
    \begin{align*}
        Lu = 0 &\text{ in }\Omega \\
        \pdv{u}{\nu} = 0 &\text{ on }\partial\Omega
    \end{align*}
    Let \(M = \sup_{\overline\Omega} u \ge 0\). We can assume this since we can just take \(-u\) instead. By the strong maximum principle, if \(u \ne M\) on \(\Omega\), then there exists \(y \in \partial\Omega\), such that \(u(y) = M\), and \(u(x) < u(y)\) for all \(x \in \Omega\). But by Hopf,
    \[{\pdv{u}{\nu}}(y) = 0\]
    Contradiction.
\end{proof}

\begin{remark}
    This says that the trivial Neumann problem (i.e. with zero data) has solutions which are constants. Now
    \[LM = 0\]
    but if \(LM = M c(x)\) for all \(x\), and so if \(c\) is not identically zero, then \(M = 0\). This constant is only non-zero when \(c = 0\)\footnote{But this is obvious, since \(Lu\) in this case only involves derivatives of \(u\).}.
\end{remark}

What happens when for non-zero right hand side? The following will be critical for Schauder theory.

\begin{theorem}
    [maximum principle a priori estimate] Suppose \(\Omega \sse \R^n\) is a bounded domain, \(L\) elliptic, \(c \le 0\), and \(\beta = \abs{b}/\lambda\) is \(L^\infty\). Let \(u \in C^2(\Omega) \cap C^0(\overline\Omega)\) and \(f : \Omega \to \R\), then
    \begin{enumerate}[(i)]
        \item if \(Lu \ge f\), then
        \[\sup_\Omega u \le \sup_{\partial\Omega} u^+ + C\sup_\Omega\left(\frac{\abs{f}}{\lambda}\right)\]
        \item if \(Lu = f\), then
        \[\sup_\Omega \abs{u} \le \sup_{\partial\Omega}\abs{u} + C\sup_\Omega\left(\frac{\abs{f}}{\lambda}\right)\]
    \end{enumerate}
    where \(C\) is a constant which only depends on \(\beta\) and \(\diam(\Omega)\).
\end{theorem}

\begin{proof}
    Set \(d = \diam(\Omega) = \sup_{x, y \in \Omega}\abs{x - y}\). As \(\Omega\) is bounded, we can contain
    \[\Omega \sse \left\{x \mid a \le x_1 \le a + d\right\}\]
    for some \(a \in \R\). Without loss of generality, \(a = 0\). As before, we will construct subsolutions, and use the weak maximum principle. Let
    \[v(x) = \sup_{\partial\Omega}u^+ + (e^{\alpha d} - e^{\alpha x_1})\sup_{\Omega}\frac{\abs{f}}{\lambda}\]
    where \(\alpha\) is to be determined.

    We can compute \(Lv\):
    \[(a^{ij}\partial_i\partial_j + b^i\partial_i)e^{\alpha x_1} = e^{\alpha x_1}(a^{11}\alpha^2 + b^1\alpha) \ge e^{\alpha x_1}\lambda(\alpha^2 + \beta \alpha) \ge \lambda\]if we take \(\alpha = \beta + 1\). Hence
    \[Lv \le cv - \lambda \sup\frac{\abs{f}}{\lambda} \le -\lambda \sup_\Omega\frac{\abs{f}}{\lambda}\]
    as \(c \le 0, v \ge 0\). Then:
    \begin{enumerate}[(i)]
        \item if \(Lu \ge f\), then
        \[L(u - v)\ge f + \lambda \sup\frac{\abs{f}}{\lambda} \ge 0\]
    \end{enumerate}
    
    \lec

    Note that \(u \le u^+\), and \(u \vert_{\partial\Omega} \le v\vert_{\partial\Omega}\) from the definition of \(v\). So by the weak maximum principle, we have \(u \le v\) in \(\Omega\), and
    \[\sup_\Omega \le u \le \sup_\Omega v \le \sup_{\partial\Omega} u^+ + C\sup_\Omega \frac{\abs{f}}{\lambda}\]
    where
    \[C = \sup_\Omega\left(e^{(\beta + 1)d} - e^{(\beta + 1)x_1}\right)\]
    For (ii), if \(Lu = f\), apply (i) to \(-u\).
\end{proof}

\subsection{H\"older spaces}

Fix \(\Omega \sse \R^n\) open, and let \(\alpha \in \Ioc{0, 1}\).

\begin{definition}
    we say \(u : \Omega \to \R\) is \emph{uniformly H\"older continuous with exponent \(\alpha\)}, or \emph{uniformly \(\alpha\)-H\"older continuous}, if
    \[[u]_{\alpha, \Omega} = \sup_{x, y \in \Omega}\frac{\abs{u(x) - u(y)}}{\abs{x - y}^\alpha} < \infty\]
    This is the \emph{H\"older seminorm}.
\end{definition}

If \(\alpha = 1\), this says that \(u\) is uniformly Lipschitz. If \(\alpha > 1\), then \(u\) is constant, by the mean value theorem.

\begin{definition}
    We say that \(u\) is \emph{locally \(\alpha\)-H\"older continuous} if for all \(K \Subset \Omega\), \(u\vert_K : K \to \R\) is uniformly \(\alpha\)-H\"older continuous.
\end{definition}

Let \(k \in \N \cup \{\infty\}\). Recall for a multi-index \(\beta \in \N^n\), we have
\[\abs{\beta} = \sum_i \beta_i\]
and
\[C^k(\Omega) = \{u : \Omega \to \R \mid \DD^\beta u\text{ exists and is continuous for all }\abs{\beta} \le k\}\]

\begin{definition}
    The \emph{H\"older space} \(C^{k, \alpha}(\Omega)\) is the space
    \[C^{k, \alpha}(\Omega) = \{u \in C^k(\Omega) \mid \DD^\beta u \text{ is locally \(\alpha\)-H\"older continuous for all }\abs{\beta} \le k\}\]
    and
    \[C^{k, \alpha}(\overline\Omega) = \{u \in C^k(\Omega) \mid \DD^\beta u \text{ is uniformly \(\alpha\)-H\"older continuous for all }\abs{\beta} \le k\}\] 
\end{definition}

For \(\alpha \in (0, 1)\), we will write
\[C^\alpha(\Omega) = C^{0, \alpha}(\Omega)\]
and
\[C^\alpha(\overline\Omega) = C^{0, \alpha}(\overline\Omega)\]

Moreover, we have
\[C^{k, 0}(\Omega) = C^k(\Omega) \qqtext{and} C^{k, 0}(\overline\Omega) = C^k(\overline\Omega)\]

\begin{remark}
    On the other hand, note that \(C^{k+1}(\Omega) \ne C^{k, 1}(\Omega)\), since Lipschitz continuity does not imply \(C^1\). On the other hand, Lipschitz functions are differentiable almost everywhere.
\end{remark}

Finally, define
\[C^{k, \alpha}_0(\Omega)=  C^{k, \alpha}_c(\Omega) = \{u \in C^{k, \alpha}(\Omega) \mid \Supp(u) \text{ compact}\}\]

To define norms on these spaces: for \(k \in \N, u \in C^k(\overline\Omega)\), define
\[[u]_{k, \Omega} = [\DD^k u]_{0, \Omega} = \sup_{\abs{\beta} = k}[\DD^\beta u]_{0, \Omega} = \sup_{\abs{\beta} = k}\sup_{x \in \Omega}\abs{\DD^\beta u(x)}\]

For \(u \in C^{k, \alpha}(\overline \Omega)\), define
\[[u]_{k, \alpha; \Omega} = [\DD^ku]_{\alpha, \Omega} = \sup_{\abs{\beta} = k}[\DD^ku]_{\alpha, \Omega}\]

Note that these are seminorms. To get norms, set
\[\norm{u}_{C^k(\Omega)} = \abs{u}_{k, \Omega} = \abs{u}_{k, 0, \Omega} = \sum_{j=0}^k\abs{D^j u}_{0, \Omega}\]
and
\[\norm{u}_{C^{k, \alpha}(\overline\Omega)} = \abs{u}_{k, \alpha, \Omega} = \abs{u}_{k, \Omega} + [\DD^k u]_{\alpha, \Omega}\]

With these norms, \(C^k\) and \(C^{k, \alpha}\) become Banach spaces. Since we will be using sequences, it is important to understand the compactness properties.

\begin{theorem}
    [Arzel\`a-Ascoli for H\"older spaces] Let \(\Omega \sse \R^n\) be a domain, \(k \in \N\), and \(\alpha \in \Ioc{0, 1}\). If \((u_j)_j\) is a sequence in \(C^{k, \alpha}(\Omega)\) such that
    \[\sup_j \abs{u_j}_{k, \alpha, \Omega'} < \infty\]
    for all \(\Omega' \Subset \Omega\), then there exists \(u \in C^{k, \alpha}(\Omega)\), and a subsequence \((u_{j'})_{j'}\) such that \(u_{j'} \to u\) in \(C^k(\overline{\Omega'})\).
\end{theorem}

\begin{remark}
    Nothing is said about convergence in \(C^{k, \alpha}\).
\end{remark}

\begin{proof}
    Examples sheet 2.
\end{proof}

We have two more ingredients before starting Schauder theory. The first is interpolation.

If we have Banach spaces \(X \Subset Y \sse Z\), then we can bound the norm in \(Y\) by \(X\) and \(Z\) norms. Interpolation is the exchange of sizes of the \(X\) and \(Z\) norms. Here, we have
\[C^{k, \alpha}(\overline\Omega) \Subset C^k(\overline\Omega) \sse C^k(\overline\Omega)\]

\begin{theorem}
    [interpolation inequality for H\"older spaces] Let \(\epsilon > 0\), \(\ell \in \N\), \(\alpha \in \Ioc{0, 1}\). Then there exists \(C =  C(n, \ell, \alpha, \epsilon)\), such that if \(u \in C^{\ell, \alpha}(\overline{B_R(x_0)})\), then
    \[R^k\abs{D^k u}_{0, B_R(x_0)} \le \epsilon R^{\ell + \alpha}[D^\ell u]_{\alpha, B_R(x_0)} + C\abs{u}_{0, B_R(x_0)}\]
    for all \(0 \le k \le \ell\).
\end{theorem}

\begin{proof}
    [Sketch proof] For details, see examples sheet 2. By rescaling and shifting, it suffices to consider the case \(R = 1\). Then argue by contradiction using Arzel\`a-Ascoli.
\end{proof}

The second ingredient is the following lemma.

\begin{theorem}
    [Simon's absorbing lemma] Let \(B_R(x) \sse \R^n\) be fixed, and let \(S\) be a nonnegative, subadditive function on the collection of sub-balls of \(B_R(x)\). That is, if
    \[B_\rho(y) \sse \bigcup_{j=1}^N B_{\rho_j}(y_j) \sse B_R(x)\]
    then
    \[S(B_\rho(y)) \le \sum_{j=1}^N S(B_{\rho_j}(y_j))\]

    Let \(\lambda \ge 0, \theta \in (0, 1)\), then there exists \(\delta = \delta(n, \lambda, \theta) \in (0, 1)\) such that: Suppose that for \emph{all} balls \(B_\rho(y) \sse B_R(x)\), we have
    \[\rho^\lambda S(B_{\theta\rho}(y)) \le \delta\rho^\lambda S(B_\rho(y)) + \gamma\]
    for some fixed \(\gamma\). Then
    \[R^\lambda S(B_{\theta R}(x)) \le C\gamma\]
    for some \(C = C(n, \theta, \lambda)\).
\end{theorem}

\begin{remark}
    This says that if there exists a local bound on \(S\), then we can ``absorb'' the \(S\)-term on the right hand side to get a global bound.
\end{remark}

\lec

\begin{proof}
    [Proof of Simon's absorbing lemma]

    Let
    \[Q = \sup_{B_\rho(y) \sse B_R(x)}\rho^\lambda S(B_{\theta\rho}(y))\]
    Recall we have that
    \[\rho^\lambda S(B_{\theta\rho}(y)) \le \delta\rho^\lambda S(B_\rho(y)) + \gamma\]
    By subadditivity of \(S\), we have that
    \[Q \le R^\lambda S(B_R(x)) < \infty\]

    Fix any \(B_\rho(y) \sse B_R(x)\). Cover \(B_{\theta\rho}(y)\) by a collection of balls
    \[\{B_{(1-\theta)\theta^2\rho}(z_j)\}_{j=1}^N\]
    with \(N \le C(\theta, n)\), which is independent of \(\rho\) and \(y\). Moreover, \(z_j \in B_{\theta\rho}(y)\). To do this: Choose a maximal pairwise disjoint collection of balls
    \[\{B_{(1-\theta)\theta^2\rho/2}(z_j)\}_{j=1}^N\]
    where \(z_j \in B_{\theta\rho(y)}\). We claim that these \(z_j\)'s work. If not, then there exists
    \[z \in B_{\theta\rho}(y) \setminus \bigcup_{i=1}^N B_{(1 - \theta)\theta^2\rho}(z_j)\]
    and so \(d(z, z_j) \ge (1 - \theta)\theta^2\rho\) for all \(j\). In particular,
    \[B_{(1-\theta)\theta^2\rho/2}(z) \cap B_{(1-\theta)\theta^2\rho/2}(z_j) = \emptyset\]
    This contradicts maximality.

    For the bound on \(N\), note that from considering radii,
    \[\bigcup_{j=1}^N B_{(1-\theta)\theta^2\rho/2}(z_j) \sse B_{(1-\theta)\theta^2\rho/2 + \theta\rho}(y) \sse B_\rho(y)\]
    Since the balls on the left hand side are disjoint, there exists a volume bound
    \[N \omega_n\left(\frac{(1 - \theta)\theta^2\rho}{2}\right)^n \le \omega_n\left(\frac{(1 - \theta)\theta^2\rho}{2} + \theta\rho\right)\]
    Which is independent of \(\rho\) and \(y\).

    To conclude, by subadditivity,
    \begin{align*}
        \rho^\lambda S(B_{\theta\rho}(y)) &\le \rho^\lambda \sum_{j=1}^N S(B_{(1 - \theta)\theta^2\rho}(z_j)) \\
        &\le \left((1 - \theta)\theta\right)^{-\lambda}\sum_{j=1}^N\left(\delta\left((1 - \theta)\theta\rho\right)^\lambda S(B_{(1 - \theta)\theta \rho}(z_j)) + \gamma\right) \\
        &\le \delta((1 - \theta)\theta)^{-\lambda}NQ + N_\gamma((1 - \theta)\theta)^{-\lambda}
    \end{align*}
    Now taking the supremum over all \(B_{\rho(y)} \sse B_R(x)\),
    \[Q \le \delta C_1 Q + C_2 \gamma\]
    where \(C_1, C_2\) depends on \(n, \theta, \lambda\). Choose \(\delta\) sufficiently small such that \(\delta C_1 < 1\). Taking \(\delta = 1/(2C_1)\),
    \[Q \le 2C_2\gamma\]
\end{proof}

\section{Schauder theory}

\subsection{Interior Schauder estimates}

We will first prove interior estimates in the unit ball, and then extend them to more general domains. The main point is: if the coefficients of \(L\) are \(\alpha\)-H\"older continuous, then any \(C^{2, \alpha}\)-solution of \(Lu = f\) can be bounded in \(C^{2, \alpha}\) on a smaller ball by \(\abs{u}_0\) and \(f\).

\begin{theorem}
    [unit scale interior Schauder estimates]
    \label{thm:unit-scale-int-schauder}

    Let \(\alpha \in (0, 1), \beta > 0\), and suppose \(a^{ij}, b^i, c \in C^{0, \alpha}(B_1(0))\), with 
    \[\abs{a^{ij}}_{0, \alpha; B_1(0)} + \abs{b^i}_{0, \alpha; B_1(0)} + \abs{c}_{0, \alpha; B_1(0)} \le \beta\]
    Suppose \(L\) is strictly elliptic, so there exists \(\lambda > 0\) such that
    \[a^{ij}(x)\xi_i\xi_j \ge \lambda \abs{\xi}^2\]
    for all \(x \in B_1(0), \xi \in \R^n\). Then if \(u \in C^{2, \alpha}(B_1(0)) \cap C^0(\overline{B_1(0)})\), and \(f \in C^{0, \alpha}(B_1(0))\) satisfies \(Lu = f\) in \(B_1(0)\), then
    \[\abs{u}_{2, \alpha; B_{1/2}(0)} \le C\left(\abs{u}_{0; B_1(0)} + \abs{f}_{0, \alpha; B_1(0)}\right)\]
    for some constant \(C = C(n, \lambda, \alpha, \beta)\).
\end{theorem}

\begin{remark}
    \begin{itemize}
        \item We can never take \(\alpha = 0\) or \(\alpha = 1\). The theorem is false in these cases.
        \item Strict ellipticity gives a lower bound for \(\lambda\), and upper bound on \(\abs{a^{ij}}_{0, \alpha}\) gives an upper bound on \(\Lambda\). So \(\Lambda/\lambda\) is bounded, and so we have uniform ellipticity.
        \item Note that we can control two derivatives of \(u\) using no derivatives on \(u\) or \(f\).
        \item 
        We will in fact strengthen this to
        \[\abs{u}_{2, \alpha; B_\theta(0)} \le C\left(\abs{u}_{0; B_1(0)} + \abs{f}_{0, \alpha; B_1(0)}\right)\]
        for all \(\theta \in (0, 1)\), \(C = C(n, \lambda, \alpha, \beta, \theta)\). 
        \item We may no assumptions, and state no conclusions, about the regularity on the boundary.
        \item The Schauder estimate gives a compactness property for the space of solutions to \(Lu = f\).
        If \((u_k) \sse C^{2, \alpha}(B_1(0)) \cap C^0(\overline{B_1(0)})\) solves \(Lu_k = f\) in \(B_1(0)\), and
        \[\gamma = \sup_k \sup_{B_1(0)}\abs{u_k} < \infty\]
        Then 
        \[\abs{u_k}_{2, \alpha; B_\theta(0)} \le C(\gamma, n, \theta, \beta, \lambda, f)\]
        and so by Arzela-Ascoli, there exists a subsequence \((u_{k'})\), with \(u \in C^{2, \alpha}(B_1(0))\), such that \(u_{k'} \to U_k\) in \(C^2(B_{\theta}(0))\) for all \(\theta \in (0, 1)\). Passing to the limit, \(Lu = f\).
    \end{itemize}
\end{remark}

\lec

\begin{proof}
    Omitted.
    % We will write \(B_r = B_r(0)\). Working in a slightly smaller ball, we can assume without loss of generality that \(\abs{u}_{2; B_1} < \infty\). The proof will consist of three steps:
    % \begin{enumerate}
    %     \item reduction,
    %     \item contradiction step,
    %     \item simplified PDE step
    % \end{enumerate}

    % \textbf{Step 1:}

    % \begin{claim}
    %     It suffices to prove: For any given \(\delta \in (0, 1)\), there exists \(C > 0\) such that
    %     \begin{equation}
    %         \label{eq:unit-scale-int-schauder-claim}
    %         [D^2u]_{\alpha; B_{1/2}} \le \delta[D^2u]_{\alpha; B_1} + C\left(\abs{u}_{2; B_1} + \abs{f}_{0, \alpha; B_1}\right)
    %     \end{equation}
    % \end{claim}

    % \begin{proof}
    %     Suppose the above holds. Then by the interpolation inequality, we obtain
    %     \begin{equation}
    %         \label{eq:unit-scale-int-schauder-claim-2}
    %         [D^2u]_{\alpha; B_{1/2}} \le 2\delta[D^2u]_{\alpha; B_1} + C\left(\abs{u}_{0; B_1} + \abs{f}_{0, \alpha; B_1}\right) 
    %     \end{equation}
    %     The strategy for this is to take a sub-ball \(B_\rho(y) \sse B_1\), and shift and scale. Define \(\widetilde u(x) = u(y + \rho x)\). Then \(\widetilde u\) will satisfy a new PDE and a new inequality.
        
    %     \lec

    %     In particular, in \(B_1\),
    %     \[\widetilde a^{ij}\partial^2_{ij}\widetilde u + \widetilde b^i \partial_i \widetilde u + \widetilde c\widetilde u = \widetilde f\]
    %     where
    %     \begin{align*}
    %         \widetilde a^{ij}(x) &= a^{ij}(y + \rho x) \\
    %         \widetilde b^i(x) &= \rho b^i(y + \rho x) \\
    %         \widetilde c(x) &= \rho^2c(y + \rho x) \\
    %         \widetilde f(x) &= \rho^2 f(y + \rho x)
    %     \end{align*}

    %     Moreover,
    %     \[\abs{\widetilde a^{ij}}_{0, \alpha; B_1} + \abs{\widetilde b^i}_{0, \alpha; B_1} + \abs{\widetilde c}_{0, \alpha; B_1} \le \abs{a^{ij}}_{0, B_\rho(y)} + \rho^\alpha[a^{ij}]_{}\]
    % \end{proof}
\end{proof}

\lec

We now give some corollaries.

\begin{corollary}
    [scale invariant Schauder estimate] 
    \label{cor:scale-invariant-schauder-estimate}
    Suppose \(B_R(x_0) \sse \R^n\), and \(a^{ij}, b^i, c \in C^{0, \alpha}(B_R(x_0))\), with
    \[a^{ij}\xi_i\xi_j \ge \lambda\abs{\xi}^2\]
    for some \(\lambda > 0\). Suppose also that
    \[\abs{a^{ij}}_{0, B_R(x_0)} + R^\alpha[a^{ij}]_{\alpha, B_R(x_0)} + R\left(\abs{b^i}_{0, B_R(x_0)} + R^\alpha[b^i]_{\alpha, B_R(x_0)}\right) + R^2\left(\abs{c}_{0, B_R(x_0)} + R^\alpha[c]_{\alpha, B_R(x_0)}\right) \le \beta\]
    for some \(\beta > 0\). Suppose \(u \in C^{2, \alpha}(B_R(x_0)) \cap C^0(\overline{B_R(x_0)})\), satisfying \(Lu = f \in C^{0, \alpha}(B_R(x_9))\). Then
    \[\abs{u}'_{2, \alpha; B_{R/2}(x_0)} \le C\left(\abs{u}_{0, B_R(x_0)} + R^2\abs{f}_{0, B_R(x_0)} + R^{2 + \alpha}[f]_{\alpha, B_R(x_0)}\right)\]
    where
    \[\abs{u}'_{k, \alpha; B_\rho(y)} = \sum_{j=0}^k \rho^j\abs{D^ju}_{0, B_\rho(y)} + \rho^{k + \alpha}[D^ku]_{\alpha, B_\rho(y)}\]
    and \(C = C(n, \lambda, \alpha, \beta)\) is independent of \(u\) and of \(R\).
\end{corollary}

\begin{proof}
    Apply \cref{thm:unit-scale-int-schauder} with \(x \to x_0 + R x\).
\end{proof}

\begin{corollary}
    [interior Schauder estimates in general domains]
    Let \(\alpha \in (0, 1), \Omega \sse \R^n\) open and bounded. Suppose we have \(a^{ij}, b^i, c \in C^{0, \alpha}(\Omega)\), with
    \[\abs{a^{ij}}_{0, \alpha; \Omega} + \abs{b^i}_{0, \alpha; \Omega} + \abs{c}_{0, \alpha; \Omega} \le \beta\]
    with
    \[a^{ij}(x)\xi_i\xi_j \ge \lambda \abs{\xi}^2\]
    for some \(\lambda > 0\). Suppose \(u \in C^{2, \alpha}(\Omega) \cap C^0(\overline\Omega)\) satisfies \(Lu = f \in C^{0, \alpha}(\Omega)\). Then for all open \(\widetilde \Omega \Subset \Omega\),
    \[\abs{u}_{2, \alpha; \widetilde \Omega} \le C\left(\abs{u}_{0, \Omega} + \abs{f}_{0, \alpha; \Omega}\right)\]
    where \(C = C(n, \alpha, \beta, \lambda, d(\widetilde\Omega, \partial\Omega))\).
\end{corollary}

\begin{proof}
    Let 
    \[d = d(\widetilde\Omega, \partial\Omega) = \sup\{r > 0 \mid \left(\widetilde\Omega\right)_r \sse \Omega\}\]
    where
    \[(\widetilde\Omega)_r = \bigcup_{x \in \widetilde\Omega}B_r(x)\]
    is the \(r\)-neighbourhood of \(\widetilde\Omega\). Then for all \(x \in \widetilde \Omega\), we have that \(B_d(x) \sse \Omega\), and so
    \[\abs{a^{ij}}'_{0, \alpha; B_d(x)} + d\abs{b^i}'_{0, \alpha; B_d(x)} + d^2\abs{c}'_{0, \alpha; B_d(x)} \le C(d)\beta\]
    Then by \cref{cor:scale-invariant-schauder-estimate}, we have an estimate
    \begin{equation}
        \begin{split}
            \abs{u}_{0; B_{d/2}(x)} + d\abs{\DD u}_{0, B_{d/2}(x)} + d^2 \abs{\DD^2 u}_{0, B_{d/2}(x)} + d^{2 + \alpha}[\DD^2u]_{\alpha, B_{d/2}(x)} &\le C\left(\abs{u}_{0, B_d(x)} + d^2\abs{f}_{0; B_d(x)} + d^{2 + \alpha}[f]_{\alpha, B_d(x)}\right) \\
            &\le C(\abs{u}_{0; \Omega} + \abs{f}_{0, \alpha; \Omega})
        \end{split}
        \tag{a}
    \end{equation}
    where \(C = C(n, \lambda, \alpha, \beta, d)\). In particular,
    \[\abs{u(x)} + \abs{\DD u(x)} + \abs{\DD^2 u(x)} \le C\left(\abs{u}_{0, \Omega} + \abs{f}_{0, \alpha; \Omega}\right)\]
    for all \(x \in \widetilde\Omega\). So
    \[\abs{u}_{2, \Omega} \le C\left(\abs{u}_{0; \Omega} + \abs{f}_{0, \alpha; \Omega}\right) \tag{b}\]
    But also, by (a),
    \[\sup_{x, y \in \widetilde\Omega, \abs{x - y} < d/2}\frac{\abs{\DD^2 u(x) - \DD^2 u(y)}}{\abs{x - y}^\alpha} \le C\left(\abs{u}_{0; \Omega} + \abs{f}_{0, \alpha; \Omega}\right)\]
    On the other hand, of \(\abs{x - y} \ge d/2\), then
    \[\frac{\abs{\DD^2u(x) - \DD^2 u(y)}}{\abs{x - y}^\alpha} \le \left(\frac{d}{2}\right)^{-\alpha}\abs{u}_{2, \widetilde\Omega} \le \left(\abs{u}_{0; \Omega} + \abs{f}_{0, \alpha; \Omega}\right)\]
    by (b). Hence
    \[[\DD^2 u]_{\alpha; \widetilde\Omega} \le \left(\abs{u}_{0; \Omega} + \abs{f}_{0, \alpha; \Omega}\right)\]
    and so combining this with (b), we get the required result.
\end{proof}

\subsection{Boundary Schauder estimates}

Write
\begin{align*}
    \R^n_+ &= \left\{(x', x^n) \mid x' \in \R^{n-1}, x^n > 0\right\} \\
    \R^n_- &= \left\{(x', x^n) \mid x' \in \R^{n-1}, x^n < 0\right\} \\
    B_R^\pm(y) &= B_R(y) \cap \R^n_\pm \\
    B_R^\pm &= B_R^\pm(0) \\
    S_R(y) &= B_R(y) \cap \{x^n = 0\} \\
\end{align*}

\begin{theorem}
    [boundary Schauder estimates in a unit ball] As before, \(0 < \alpha < 1\), \(a^{ij}, b^i, c \in C^{0, \alpha}(B_1^+)\), and
    \[\abs{a^{ij}}_{0, \alpha; B_1^+} + \abs{b^i}_{0, \alpha, B_1^+} + \abs{c}_{0, \alpha; B_1^+} \le \beta\]
    and
    \[a^{ij}(x)\xi_i\xi_j \ge \lambda\abs{\xi}^2\]
    for all \(x \in B_1^+\). Suppose \(u \in C^{2, \alpha}(B_1^+)\) satisfies
    \[\begin{cases}
        Lu = f &\text{in }B_1^+ \\
        u = 0 &\text{on }S_1
    \end{cases}\]
    where \(f \in C^{0, \alpha}\) and \(\varphi \in C^{2, \alpha}\). Then
    \[\abs{u}_{2, \alpha, B_{1/2}^+} \le C\left(\abs{u}_{0; B_1^+} + \abs{f}_{0, \alpha; B_1^+} + \abs{\varphi}_{2, \alpha; B_1^+}\right)\]
\end{theorem}

\begin{proof}
    By considering \(v = u - \varphi\), suffices to consider the case when \(\varphi = 0\), since \(L\varphi \in C^{0, \alpha}(B_1^+)\).

    The rest of the proof is as in \cref{thm:unit-scale-int-schauder}, which we will omit.
\end{proof}

\begin{proposition}
    [reflection principle for harmonic functions] Let \(\Omega^+\) be an open subset of \(\R^n_+\), and let \(T = \partial\Omega^+ \cap \{x^n = 0\}\). Let \(\Omega^-\) be the reflection of \(\Omega^+\) in \(\{x^n = 0\}\). Let \(v \in C^2(\Omega^+) \cap C^0(\Omega^+ \cup T)\), and let \(\overline{v}\) be the odd reflection of \(v\) in \(T\). That is, \(\overline{v} : \Omega^+ \cup T \to \Omega^- \to \R\), where
    \[\overline v(x', x^n) = \begin{cases}
        v(x', x^n) & (x', x^n) \in \Omega^+ \cup T \\
        -v(x', -x^n) & (x', x^n) \in \Omega^-
    \end{cases}\]
    Then if \(\Delta v = 0\) in \(\Omega^+\) and \(v\vert_T = 0\), then \(\overline{v} \in C^2(\Int(\Omega^+ \cup T \cup \Omega^-))\) and \(\Delta \overline{v} = 0\).
\end{proposition}

\begin{proof}
    Use the mean value property. See examples sheet 2.
\end{proof}

\begin{remark}
    This is trivial if \(T = \emptyset\), since \(\Omega^+\) and \(\Omega^-\) are disjoint. The important part of this theorem is that \(\overline v\) is \(C^2\) across \(T\). 
\end{remark}

\begin{proposition}
    [boundary absorbing lemma] Given \(\theta \in (0, 1), \mu \in \R\), then there exists \(\delta = \delta(n, \theta, \mu)\), and \(C = C(n, \theta, \mu)\) such that:

    Fix \(R > 0\), and let
    \begin{align*}
        \mcB &= \left\{B_\rho(y) \mid B_\rho(y) \sse B_R^+(0)\right\} \\
        \mcB^+ &= \left\{B^+_\rho(y) \mid y^n = 0, B^+_\rho(y) \sse B_R^+(0)\right\}
    \end{align*}
    Suppose \(S : \mcB \cup \mcB^+ \to \R\) is a nonnegative subadditve function, such that
    \[\rho\mu S(B_{\theta\rho}^+(y)) \le \delta\rho\mu S(B_\rho^+(y)) + \gamma\]
    for all \(B_\rho^+(y) \in \mcB^+\), and
    \[\rho\mu S(B_{\theta\rho}(y)) \le \gamma\]
    for \(B_\rho(y) \in \mcB\). Then
    \[R^\mu S(B_{\theta R}^+(0)) \le C\gamma\]
\end{proposition}

\begin{proof}
    Examples sheet 2.
\end{proof}

\lec

Let (H) denote the following hypothesis: ``Suppose \(\Omega \sse \R^n\) is a bounded domain, and \(0 < \alpha < 1\). Suppose \(a^{ij}, b^i, c \in C^{0, \alpha}(\Omega)\), with
\[\abs{a^{ij}}_{0, \alpha; \Omega} + \abs{b^i}_{0, \alpha; \Omega} + \abs{c}_{0, \alpha; \Omega} \le \beta\]
and suppose there exists \(\lambda > 0\) such that
\[a^{ij}(x)\xi_i\xi_j \ge \lambda\abs{\xi}^2\]''

\begin{theorem}
    [boundary Schauder estimates in general domains]

    Suppose (H) holds, ad \(\Omega\) is a \(C^{2, \alpha}\) domain. Then there exists \(\epsilon = \epsilon(\Omega) > 0\) such that if \(u \in C^{2, \alpha}(\overline\Omega)\), \(f \in C^{0, \alpha}(\Omega), \varphi \in C^{2, \alpha}(\overline\Omega)\) such that
    \[\begin{cases}
        Lu = f &\text{in }\Omega \\
        u = \varphi&\text{on }\partial\Omega
    \end{cases}\]
    Then for all \(x \in \partial\Omega\),
    \[\abs{u}_{2, \alpha; B_\epsilon(x) \cap \Omega} \le C\left(\abs{u}_{0; \Omega} + \abs{f}_{0, \alpha; \Omega} + \abs{\varphi}_{2, \alpha; \Omega}\right)\]
\end{theorem}

\begin{remark}
    We need \(\Omega\) to be \(C^{2, \alpha}\) for \(u\) to be \(C^{2, \alpha}\) on \(\partial\Omega\).
\end{remark}

\begin{proof}
    Omitted.
    % Choose \(z \in \partial\Omega\). By definition, there eixsts \(R > 0\) and a \(C^{2, \alpha}\) diffeomorphism \(\psi : B_R(z) \to D\), sending \(B_R(z) \cap \Omega\) to \(y_n > 0\), and \(B_R(z) \cap \partial\Omega\) to \(y_n = 0\).

    % Let \(x = (x^1, \dots, x^n)\) be coordinates in \(\Omega\), and let \(y = (y^1, \dots, y^n)\) be the corresponding coordinates in \(D\). Let
    % \[\widetilde u(y) = u(\psi^{-1}(y))\]
    % be the pullback of \(u\) along \(\psi^{-1}\). Then 
    % \[\widetilde u\vert_{\{y^n = 0\} \cap D} = \left(\varphi \circ \psi^{-1}\right)\vert_{\{y^n = 0\} \cap D} = \widetilde \varphi\]

    % To apply the unit scale boundary Schauder estimates, we need to find a PDE satisfied by \(\widetilde u\), and show that it satisfies the hypothesis. Note that \(u(x) = \widetilde u(\psi(x))\), and so
    % \[\partial_{x^i}u = \partial_{y^k}\widetilde u \pdv{\psi^k}{x^i}\]
    % and
    % \[\partial^2_{x^ix^j}u = \partial^2_{y^ky^\ell}\widetilde u \pdv{\psi^k}{x^i}\pdv{\psi^\ell}{x^j} + \partial_{y^k}\widetilde u + \pdv[2]{\psi^k}{x^i}{x^j}\]
    % Hence we can find the coefficients of the new PDE explicitly, which is
    % \[A^{ij} \partial^2_{y^iy^j}\widetilde u + B^i\partial_{y^i}\widetilde u + C\widetilde u = \widetilde f\]
    % where for example,
    % \[A^{ij} = a^{k\ell}\pdv{\psi^i}{x^k}\pdv{\psi^j}{x^\ell}\]
\end{proof}

\subsection{Global Schauder estimates}

We can combine the interior and boundary estimates to get:

\begin{theorem}
    [global Schauder estimates]

    Suppose (H) holds, and suppose that \(\Omega\) is a \(C^{2, \alpha}\) domain. If \(u \in C^{2, \alpha}(\Omega), f \in C^{0, \alpha}(\Omega), \varphi \in C^{2, \alpha}(\Omega)\) is such that
    \[\begin{cases}
        Lu = f &\text{in }\Omega \\
        u = \varphi &\text{on }\partial\Omega
    \end{cases}\]
    Then
    \[\abs{u}_{2, \alpha; \Omega} \le C\left(\abs{u}_{0, \Omega} + \abs{f}_{0, \alpha; \Omega} + \abs{\varphi}_{2, \alpha; \Omega}\right)\]
    where \(C = C(n, \lambda, \alpha, \beta, \Omega)\).
\end{theorem}

\end{document}
